\section[Información]{Información del Proyecto}
\normalsize{ \indent
Sobre el proyecto se obtiene la siguiente información:
}
\begin{tabbing}
  \hspace{1cm}\= \textbf{Nombre}\qquad \qquad \qquad
  \quad \quad \= SAES(Sistema de Administración para
  \\[5pt] \> \> \qquad \quad Escuelas Secundarias) \\[5pt] \>
  \textbf{Tipo} \> Con cliente final\\[5pt] \>
  \textbf{Fecha de preparación} \>
  9 de septiembre de 2022\\[5pt] \>
  \textbf{Cliente} \> Personal del \acrlong{ilc}
  \\[5pt] \> \textbf{Patrocinador} \> \acrlong{ilc}
  \\[5pt] \> \textbf{Líder} \> Rodrigo Lionel Antoniak
\end{tabbing}
\section[Resumen]{Resumen Ejecutivo}
\normalsize{ \indent
Para este proyecto, se pretende desarrollar un
software donde haya una base de datos para almacenar
legajos, horarios, calificaciones y asistencias;
donde tenga un dominio en un sitio web, para así
poder portar el manejo de perfiles en distintas
partes. No ha de tener muchas dificultades respecto
a la competencia, ya que conserva su simplicidad en
el formato que se presenta.
}
\newline
\normalsize{ \indent
Sus puntos principales son el cumplimiento a las
necesidades del colegio, con un formato que favorece
al uso de recursos humanos que ya se poseen y poco
entrenamiento; teniendo que invertir principalmente en
la programación y la base de datos. La tarea que abarque
la totalidad del precio de la solución es la instalación
de software, debido a que se utiliza la licencia GPL v2.
}
\newline
\normalsize{ \indent
Como el diseño es bastante personalizado por el
\acrlong{ilc}, se pretende que el proyecto se realice una
única vez; sin llevarlo a la reventa del programa, o al
menos no es el propósito. Cabe destacar que cada escuela
puede solicitar requisitos distintos, por lo que la
variedad de soluciones es alta para realizar una de
propósito general.
}
\newline
\normalsize{ \indent
Visto y considerando que existe información suficiente
sobre el contexto, podrá llevarse a cabo este proyecto
(o al menos en el alcance pretendido actualmente);
después, será necesario continuar con las buenas prácticas
para completar todo el sistema e integrarlo como espera
el colegio secundario.
}
\section[Antecedentes]{Antecedentes del Proyecto}
\normalsize{ \indent
Primero, los factores que justifican la realización del
proyecto son la ausencia de digitalización de documentos,
el medio de manejo actual de información (a través de Excel)
y el fracaso de personal informático (contratado por freelance)
para realizar correctamente un proyecto funcional efectivo.
De ahí, el estudio de factibilidad es impulsado por las posibles
opciones que existan, la abertura del cliente a distintas
propuestas y la necesidad de decidir si comprar o desarrollar
software.
}
\section[Contexto]{El Proyecto y su Contexto}
\normalsize{ \indent
Una escuela de nivel medio desea digitalizar la generación
de formularios y completado de legajo. Los docentes y alumnos
pueden completar los datos de sus legajos desde una plataforma,
donde se conservará la información en una base de datos; así
posteriormente se genera el documento y es firmado
presencialmente. El único documento excepcional de esto es la
ficha médica, la cual se ofrece el modelo y se trae el papel
completado por el médico.
}
\newline
\normalsize{ \indent
El acceso a los datos de los legajos será compartido,
aunque estará limitado a aquellos que tengan el rango
suficiente para acceder a la información. Los legajos
podrán observarse desde el perfil del alumno o docente,
o en una tabla que junte a los cursantes o los profesores.
Los docentes pueden clasificarse por el tipo de salario
(por hora, cargo, o ambos), mientras que los alumnos se
pueden clasificar por curso/año/división.
}
\newline
\normalsize{ \indent
Entre otros, debe disponerse de un control de horarios;
propuesto para que la Rectora pueda asignar los docentes
correspondientes a cada rol y el sistema ajuste los
horarios del mismo. Cuando se necesite cambiar los
horarios, el sistema debe controlar que se cumplan todas
las condiciones para la operación. A su vez, la opción de
un docente suplente (dado por la licencia del titular) es
fundamental en el sistema; igualmente que el cambio de
docentes en medio del año.
}
\newline
\normalsize{ \indent
Además, debe permitir registrar asistencias con los
alumnos y docentes de forma directa; así, se libera
al preceptor para que se encargue de otros trabajos.
Cada presente abarcará la condición de un alumno
con las horas contínuas de una materia determinada,
implicando que tras el cambio de materia o la descontinuidad
de una misma materia debe contabilizarse como una asistencia
aparte.
}
\newline
\normalsize{ \indent
Por otro lado, es necesario almacenar las calificaciones
entre docentes y alumnos; tanto las que reciben estos
últimos por su desenvolvimiento en el cursado, como del
profesor por la calidad de enseñanza que otorga. El primer
tipo de calificación requiere que sea trimestral, en donde
cada materia de ese lapso debe tener 3 calificaciones cada
uno. Mientras que las devoluciones de alumnos a docentes
se realiza anónimamente de forma aleatoria.
}
\newline
\normalsize{ \indent
El sistema debe poder generar notificaciones para
informar sobre el estado de cada uno de los estados del
sistema, independientemente del perfil que se involucre.
El motivo del aviso sobre documentación puede ser la falta
de los mismos para completar el legajo o que se notifique
de estar al día, a su vez que puede destacar algún tipo de
error en los datos ingresados en un formulario en específico.
Adicionalmente, los docentes recibirán un correo a través del
dominio del colegio.
}
\newline
\normalsize{ \indent
En adición, existen otros formularios que pueden
generar los miembros por temas distintos al legajo;
entonces se necesitará un canal distinto de gestión
para esos documentos. Incluso, se deberán incluir
documentos históricos; teniendo en cuenta papeles
anteriores a la creación del sistema, siguiendo un
orden lógico de ubicación (sea por fecha, alumno,
curso, etcétera).
}
\newline
\normalsize{ \indent
Por último, debe haber un apartado donde se administren
los exalumnos; teniendo sus nombres, datos de graduación
(año/título) y contacto. Idealmente, cada uno de los
egresados podría cambiar datos de su contacto; así el
colegio puede conservar contacto con ellos. Si bien,
puede evaluarse un sistema de mensajería por el software;
el colegio insiste en tener medios alternativos, con
motivos propios del ente.
}
\newline
\normalsize{ \indent
Respecto al contexto, se puede mencionar que el autor
de este trabajo final ha cursado su nivel medio en este
colegio; lo cual, facilita la comunicación con los
miembros del establecimiento. Tomando en cuenta el
diálogo en la entrevista, puede afirmarse que hubo
acuerdo en la visión de cada miembro de Coordinación
Escolar entrevistado.
}
\section[Alcance]{Alcance del Estudio de Factibilidad}
\normalsize{ \indent
El alcance pretendido para este estudio es el conocimiento
de los pasos a seguir para poder gestionar legajos,
horarios, calificaciones y asistencias. Principalmente,
se analizarán distintas opciones por los costos de
almacenamiento y dinero para elegir; considerando que debe
cumplirse la condición de poder usar el personal ya
contratado del colegio con la posible solución.
}
\section{Factibilidad Técnica}
\normalsize{ \indent
En la búsqueda de las distintas soluciones,
se han encontrado distintas herramientas que
parecen ser candidatas a ser parte del software
final. Para la investigación inicial, se enfocó
en encontrar soluciones destinadas a escuelas;
donde se incluya el manejo de legajos y horarios
docentes. Además, se recuerda que el colegio ya
posee el software de Acadeu; tras el suceso de que
el software no satisface las necesidades del
\acrlong{ilc}, se debe módulos de calificaciones y
asistencia. Gestionar tareas para alumnos debe 
revisarse para un proyecto a futuro, considerando
que la vuelta a la normalidad (frente a la situación
del COVID-19) permite la entrega presencial de
tareas y se insiste en conservar este método.
}
\newline
\normalsize{ \indent
Con todos los filtros establecidos, se destaca un
proyecto de código abierto (bajo licencia pública de
GNU); se llama Proyecto ALBA, aunque ya fue abandonado.
Consiste en la gestión de unidades educacionales, junto
con sus alumnos y docentes; destacándose que existen
módulos para gestionar legajos y horarios docentes. El
asunto es que no parece ser muy conocido y no tuvo mucha
atención, ya que todos los problemas reportados no fueron
respondidos (y quedaron abiertos). Incluso, su página web
principal ya no existe; empero su enlace por GitHub sigue
siendo accesible (como archivado), el cual es el siguiente:
\url{https://github.com/proyectoalba/alba}.
}
\newline
\normalsize{ \indent
Por lo tanto, se observa que no hay una opción íntegra
que cumpla con lo pretendido en el alcance del proyecto;
así que se ha buscado las funcionalidades por separado, en
caso que existiera una herramienta disponible. Si bien, para
la administración de horas docentes se encuentran herramientas
con propósitos laborales generales; no hay alguna que se centre
solamente en gestionar horarios, es decir que trae muchas otras
funcionalidades no deseadas. Se puede nombrar de ejemplo a
PlanningPME
(\url{https://www.planningpme.es/software-planning.htm})
y Sesame (\url{https://www.sesametime.com/}). Para la
gestión de legajos, asistencia y calificaciones; se considera
que la funcionalidad es suficientemente genérica como para
poder generar el código sencillamente y que se enfoque en
solucionar específicamente lo deseado.
}
\section{Factibilidad Económica}
\normalsize{ \indent
Este factor puede resultar esencial hasta cierto punto,
pero hay consideraciones a tomar que bajan la importancia
de este factor:
}
\begin{itemize}
  \item Este proyecto no se pretende volver a usar, ya que
  el colegio nos ha elegido en criterio a la relación de 
  alumno (por el autor de este documento); mas falta decir
  que deja de ser un proyecto cuando se vuelve a realizar.
  Otros colegios pretenderían buscar otras características,
  por lo que sería un proyecto aparte.
  \item El \acrlong{ilc} tiene predisposición a pagar precios
  elevados, siempre y cuando se cumpla con la funcionalidad
  buscada. Además, tiene conocimientos sobre los precios
  reales de software (ya se comentó antes que intentaron
  tener este proyecto resuelto antes con otras personas, pero
  estos últimos no estuvieron a la altura de las
  circunstancias).
  \item Dado que el personal se pretende conservar y que no
  se necesitará mucho entrenamiento, gracias a que se sigue
  el formato que ellos desean con los procedimientos
  adecuados para \Gls{coordesc}; los únicos elementos
  a cobrar son la base de datos, el dominio y la programación.
\end{itemize}
\ \newline
\normalsize{ \indent
Dado que no se conoce la programación requerida, resulta
imposible estimar en este tiempo el costo en sí; sin embargo,
no resulta tan relevante frente a otros proyectos con
competencia y recursos financieros muy limitados.
}
\section{Factibilidad Legal}
\normalsize{ \indent
Tomando en cuenta que se pretende hacer uso de una base de
datos para almacenar los datos sobre los docentes (contando
sólo esta etapa), debe considerarse el manejo de los datos
personales; cuya protección está normalizada en la ley 25326,
que se encuentra en el siguiente enlace:
\url{http://servicios.infoleg.gob.ar/infolegInternet/anexos/60000-64999/64790/norma.htm}.
}
\newline
\normalsize{ \indent
De ahí, hay que tener énfasis en el artículo 10; en
consecuencia, el secreto profesional de los datos personales
es obligatorio (incluso después de terminar la solución
completa).
}
\section{Factibilidad de Recursos}
\normalsize{ \indent
Desde el inicio del proyecto, se tuvo la visión de un relevo
de métodos para realizar operaciones; y no de agregar más
trabajo para hacer, por lo que se aplicó esfuerzo en
limitarse con los recursos ya existentes del colegio.
Gracias a que el personal de \Gls{coordesc} ya realiza
todas las tareas pretendidas con este sistema, aunque con
una forma poco efectiva; no se necesitará agregar personal
para este sistema, ni se afecta a ninguna operación (por lo
que no hay dependencias ni procedimientos necesarios para
describir).
}
\newline
\normalsize{ \indent
Apartando los recursos humanos ya presentes, es necesario
tener una base de datos para el almacenamiento de datos
sobre legajos y horarios (inclusive también para la
incorporación de los elementos pertenecientes a las
siguientes etapas). Para este proyecto, tanto en el
alcance actual como en su totalidad, es ideal realizar
con un dominio web; por motivos de portabilidad para los
docentes y alumnos, cuando quieran actualizar su perfil.
En caso de elegir un diseño de página web, puede
aprovecharse el dominio ya existente del colegio; el cual
es \url{https://www.ilc.edu.ar/}.
}
\section{Factibilidad de Mercado}
\normalsize{ \indent
Como se ha visto en la factibilidad técnica, existen
varias opciones en el mercado que se acercan a algunas
de las necesidades. Para la gestión de una escuela, hay
muchas opciones que se hallan buscando en Internet; sin
embargo, poseen características que resultan ser
complicadas de implementar en muchos colegios. Por
ejemplo, el seguimiento de asistencia, tomando en cuenta
que el colegio cuenta con una residencia para internados;
en especial, por tener una ubicación biométrica o GPS.
}
\newline
\normalsize{ \indent
Por otro lado, hay aplicaciones que se destinan a la
gestión de horarios; la cual administra hasta las
vacaciones que poseen cada uno de los trabajadores. El
problema es que no gestiona los legajos docentes, así
que lamentablemente no se tiene suficiente
especialización en estas alternativas. En fin, esta
solución a desarrollar va a resultar como una oferta
única en el mercado; aunque se pretende que el
\acrshort{ilc} utilice, sin buscar que otros colegios
puedan pedir tal oferta.
}
\newline
\normalsize{ \indent
Para resumir, este software está pretendido para que
los colegios puedan gestionar la administración pedagógica;
sin tener herramientas complejas de implementar en
regiones no pertenecientes al primer mundo. En principio,
el primer cliente (el colegio secundario que solicitó
nuestra solución) elige nuestro sistema porque va a cumplir
correctamente la función solicitada; sin incorporar
herramientas innecesarias para esta escuela, como otras
ofertas.
}
\newline
\normalsize{ \indent
En otras palabras, la persona que escribe este documento
tendrá una solución que está entremedio de dos tipos de
competencias. Por un lado, están los programadores a nivel
local; quienes no resultan ser suficientemente competentes
para cumplir lo que busca el colegio, según testimonio del
personal del establecimiento. El otro sector es el conjunto
de organizaciones que ofrece demasiadas herramientas para
colegios de regiones lejanas a ciudades como Córdoba.
Ofrece una oportunidad donde se tienen las características
programadas correctamente y mantiene la simplicidad del
mismo, siendo el enfoque de esta solución.
}
\section{Factibilidad Operacional}
\normalsize{ \indent
Dentro de las operaciones que se cumplen, se gestiona
la base de datos con los legajos y horarios. Para los
legajos, se completa el legajo de los docentes por dos
formas; dependiendo de quién haga la carga de los archivos
digitales, sean las imágenes o la declaración jurada
firmada. Una persona puede ser el mismo docente,
rellenando un formulario con información que se
complementa con los cargos asignados, los títulos
superiores y la información básica de perfil; para así
tener el modelo de declaración jurada digital que se
podrá imprimir y firmar para guardarse en la base de
datos. La otra opción es el personal de \Gls{coordesc},
donde éste deberá realizar la carga de la información
del formulario mencionado anteriormente; considerando
que el docente prefiere completar la declaración jurada
a mano.
}
\newline
\normalsize{ \indent
Todo lo explicado con respecto a la declaración jurada
se aplica igualmente con las imágenes, distinguiéndose
la verificación de imagen; en el caso que el docente
suba la imagen por su decisión, ya que no hay seguridad
de que la imagen corresponda al documento solicitado
(incluso, puede tratarse de una confusión de documentación;
donde un docente sube un documento en el lugar que
corresponde a otro).
}
\newline
\normalsize{ \indent
Todo lo explicado con respecto a la declaración jurada
se aplica igualmente con las imágenes, distinguiéndose
la verificación de imagen; en el caso que el docente
suba la imagen por su decisión, ya que no hay seguridad
de que la imagen corresponda al documento solicitado
(incluso, puede tratarse de una confusión de documentación;
donde un docente sube un documento en el lugar que
corresponde a otro).
}
\newline
\normalsize{ \indent
Para los horarios, la Rectora es quien asigna los
docentes que se encargan de los distintos puestos
disponibles; subdividiéndose en horas jornales y cargo,
donde un docente puede tener ambos. Así también, la
Rectora hace los \glspl{reldoc}; sea por una suplencia o
por un reemplazo de titular. De ahí, \Gls{coordesc} se
encarga de ajustar los horarios según las condiciones
necesarias para ello; tanto el ajuste inicial de cada
año como los cambios de horario interanuales.
}
\newline
\normalsize{ \indent
Al elegir un docente para cierto horario, la Rectora
tiene todos los permisos para hacer la operación y no
necesita guiarse de ninguna condición. Por el contrario,
para los horarios existe un conjunto de necesidades a
cumplir para los cambios:
}
\begin{itemize}
  \item Un horario a cambiar de un docente no puede
  sobreponerse a otro del mismo profesional, ya sea en
  el mismo colegio o externamente (dato conocido en la
  declaración jurada).
  \item Para una clase que no corresponda a \Gls{modulo}
  (práctica), no se puede tener más de tres horas
  consecutivas de cátedra.
  \item Las horas de \Gls{modulo} deben ser 7:30 a 12:00,
  9:00 a 12:00 (en ambos casos, con un recreo de 10:00
  a 10:30 y sin actividades después del mediodía), o
  14:00 a 16:30 (sin actividades en adelante).
  \item Los docentes no pueden trabajar más de 42 horas.
  Para los que trabajan fuera del colegio (quienes sólo
  tengan horas jornales), se considera el horario externo.
  \item Las horas correspondientes a un curso no se pueden
  interponer. Por lo tanto, siempre que se cambien horarios;
  resulta en una rotación, donde los horarios involucrados
  pueden ser variables (al menos son dos).
\end{itemize}
\ \newline
\normalsize{ \indent
Adicionalmente, los \Glspl{modulo} no pueden cambiarse de
horario dentro del año.
}
\section{Factibilidad de Tiempo}
\normalsize{ \indent
Considerando que la competencia no cubre las necesidades de
la escuela específicamente, no hay un riesgo aparente del
tiempo frente a otros que puedan hacer la misma tarea en
menor tiempo; porque no se halló a alguien con esos rasgos.
Sin embargo, no hay que olvidarse que uno se compromete con
la entrega del software; por lo tanto, se debe seguir los
tiempos previamente tratados con quien nos vaya a comprar el
sistema.
}
\section{Recomendaciones}
\normalsize{ \indent
Cuando se siga con otros objetivos generales que no se abarcan
en el alcance de ahora, es recomendable conservar las prácticas
actuales; así en el futuro se podrá apuntar al éxito de este
proyecto, el cual ya contempla un gran potencial.
}
\newline
\normalsize{ \indent
Ya que no hay muchas opciones en el mercado que cumplan las
funciones del alcance de forma simplificada, se consiguen los
motivos para llevar adelante el software pretendido. Además,
el acuerdo entre los desarrolladores y el cliente de la forma
final otorga confianza que es una buena idea para seguir.
}
\newline
\normalsize{ \indent
Para finalizar, es muy factible ejecutar el proyecto en
cuestión (al menos, hasta el alcance en esta primera etapa) y
debe continuarse con los procedimientos correspondientes; para
así lograr los objetivos de la solución que se creará.
}
